\documentclass[14pt, a4paper]{extarticle}
\usepackage[utf8]{inputenc}
\usepackage[english,ukrainian]{babel}
\usepackage{amsmath, amssymb, amsfonts, amsthm}
\usepackage{geometry}
\usepackage{graphicx}
\usepackage{hyperref}
\usepackage{enumitem}
\usepackage{array}
\usepackage{booktabs}
\usepackage{xcolor}
\usepackage{colortbl}
\usepackage{longtable}
\usepackage{multirow}

\usepackage{fontspec}
\setmainfont{Times New Roman}
\usepackage[T2A]{fontenc}

\geometry{top=2cm, bottom=2cm, left=0.1cm, right=2cm}

\hypersetup{
    colorlinks=true,
    linkcolor=blue,
    urlcolor=blue,
}

% Нумерація формул за розділами
\numberwithin{equation}{section}

% Стиль для позначень
\newtheorem{notation}{Позначення}[section]

\title{Повний список формул\\інформаційної технології ДІПДО\\
\large{Динамічна Ідентифікація та Параметризація Динамічних Об'єктів}}
\author{Наукове дослідження}
\date{\today}

\begin{document}
\section{МЕТОДИ ДИСТАНЦІЙНОЇ ІДЕНТИФІКАЦІЇ НА ОСНОВІ КОМП'ЮТЕРНОГО ЗОРУ ТА ГЛИБИННОГО НАВЧАННЯ}

У цьому розділі розглядаються математичні методи та моделі, що використовуються для дистанційної ідентифікації динамічних об'єктів. Основна увага приділяється методам оптичного потоку, нейромережевим архітектурам для аналізу руху, ансамблевим підходам та методам об'єднання різних джерел інформації.

\subsection{Методи оптичного потоку в задачах ідентифікації}

Оптичний потік є фундаментальним поняттям у комп'ютерному зорі, що описує розподіл видимих швидкостей руху об'єктів на зображенні. Нехай $I(x,y,t)$ — функція яскравості зображення в точці $(x,y)$ у момент часу $t$. Тоді основне рівняння оптичного потоку має вигляд \cite{horn1981determining}:

\begin{equation}
I_x u + I_y v + I_t = 0 \tag{8}
\end{equation}

де $u = \frac{dx}{dt}$, $v = \frac{dy}{dt}$ — компоненти оптичного потоку, а $I_x = \frac{\partial I}{\partial x}$, $I_y = \frac{\partial I}{\partial y}$, $I_t = \frac{\partial I}{\partial t}$ — часткові похідні яскравості.

\subsubsection{Локальний метод Лукаса-Канаде}

Метод Лукаса-Канаде \cite{lucas1981iterative} базується на припущенні, що оптичний потік є локально постійним у межах невеликого вікна розміром $N \times N$. Для пікселів вікна формується система рівнянь:

\begin{equation}
\begin{bmatrix} 
I_x(p_1) & I_y(p_1) \\ 
I_x(p_2) & I_y(p_2) \\ 
\vdots & \vdots \\ 
I_x(p_n) & I_y(p_n) 
\end{bmatrix} 
\begin{bmatrix} u \\ v \end{bmatrix} = 
-\begin{bmatrix} I_t(p_1) \\ I_t(p_2) \\ \vdots \\ I_t(p_n) \end{bmatrix} \tag{9}
\end{equation}

Розв'язок цієї системи методом найменших квадратів дає оцінку оптичного потоку:

\begin{equation}
\begin{bmatrix} u \\ v \end{bmatrix} = (A^T A)^{-1} A^T b \tag{10}
\end{equation}

де $A$ — матриця градієнтів інтенсивності, $b$ — вектор часових різниць. Перевагою методу є його обчислювальна ефективність та стійкість до шумів, однак він не здатен відновлювати потік в областях з однорідною текстурою.

\subsubsection{Глобальний метод Хорна-Шунка}

На відміну від локального підходу, метод Хорна-Шунка \cite{horn1981determining} накладає глобальну умову гладкості на поле оптичного потоку. Мінімізується наступний функціонал:

\begin{equation}
E = \iint \left[(I_x u + I_y v + I_t)^2 + \alpha^2 \left( \|\nabla u\|^2 + \|\nabla v\|^2 \right) \right] dx\,dy \tag{11}
\end{equation}

Перший доданок відповідає за виконання базового рівняння оптичного потоку, а другий — за гладкість поля. Параметр $\alpha$ регулює баланс між цими двома умовами. Розв'язок знаходиться ітераційно:

\begin{align}
u^{k+1} &= \bar{u}^k - \frac{I_x (I_x \bar{u}^k + I_y \bar{v}^k + I_t)}{\alpha^2 + I_x^2 + I_y^2} \tag{14} \\
v^{k+1} &= \bar{v}^k - \frac{I_y (I_x \bar{u}^k + I_y \bar{v}^k + I_t)}{\alpha^2 + I_x^2 + I_y^2} \tag{15}
\end{align}

де $\bar{u}^k$, $\bar{v}^k$ — усереднені значення потоку на $k$-й ітерації.

\subsubsection{Поліноміальний метод Farneback}

Метод Farneback \cite{farneback2003two} базується на поліноміальній апроксимації зображення. Локальне поліноміальне представлення має вигляд:

\begin{equation}
f(x) \approx x^T A x + b^T x + c \tag{1}
\end{equation}

Для двох послідовних кадрів припускається, що вони пов'язані зсувом $d$: $f_2(x) = f_1(x-d)$. Після розкладання та порівняння коефіцієнтів отримуємо:

\begin{align}
A_2 &= A_1 \tag{4} \\
b_2 &= b_1 - 2A_1 d \tag{5} \\
c_2 &= d^T A_1 d - b_1^T d + c_1 \tag{6}
\end{align}

Звідси вектор зміщення обчислюється як:

\begin{equation}
d = -\frac{1}{2} A_1^{-1} (b_2 - b_1) \tag{7}
\end{equation}

Цей метод забезпечує високу точність і є стійким до шумів.

\subsection{Нейромережеві моделі для аналізу руху}

Сучасні методи глибинного навчання дозволяють значно підвищити точність оцінки оптичного потоку та глибини.

\subsubsection{Модель FlowNet}

FlowNet \cite{dosovitskiy2015flownet} є однією з перших згорткових нейронних мереж для задачі оптичного потоку. Архітектура має структуру енкодер-декодер:

\begin{equation}
\Phi_{\text{encoder}}(I_1, I_2) = \text{Conv}_{1 \times 1}(\text{Conv}_{3 \times 3}(\cdots \text{Conv}_{7 \times 7}([I_1, I_2])\cdots)) \tag{16}
\end{equation}

\begin{equation}
\hat{F} = \Phi_{\text{decoder}}(\Phi_{\text{encoder}}(I_1, I_2)) \tag{17}
\end{equation}

Оптимізація мережі відбувається шляхом мінімізації функції втрат EPE (End-Point Error):

\begin{equation}
\mathcal{L}_{\text{EPE}} = \frac{1}{N} \sum_{i=1}^N \|F_{\text{pred}}(i) - F_{\text{gt}}(i)\|_2 \tag{18}
\end{equation}

\subsubsection{Модель RAFT}

RAFT (Recurrent All-pairs Field Transforms) \cite{teed2020raft} використовує кореляційний об'єм для встановлення відповідностей між кадрами:

\begin{equation}
C_{ijkl} = \frac{1}{\sqrt{HW}} \sum_{h,w} f_{ijhw} \cdot g_{klhw} \tag{19}
\end{equation}

Оновлення поля потоку виконується за допомогою GRU (Gated Recurrent Unit):

\begin{equation}
\Delta x_t, h_{t+1} = \text{GRU}([x_t, C(x_t)], h_t) \tag{21}
\end{equation}

Такий рекурентний підхід дозволяє отримувати високу точність на великих зміщеннях.

\subsubsection{Модель DETR для виявлення об'єктів}

DETR (Detection Transformer) \cite{carion2020end} використовує архітектуру трансформера для кінцево-кінцевого виявлення об'єктів. Механізм багатоголової уваги обчислюється як:

\begin{equation}
\text{MultiHead}(Q, K, V) = \text{Concat}(\text{head}_1, \ldots, \text{head}_h) W^O \tag{58}
\end{equation}

де кожна голова визначається як:

\begin{equation}
\text{head}_i = \text{Attention}(QW_i^Q, KW_i^K, VW_i^V) \tag{59}
\end{equation}

Оптимальне парування між передбаченнями та істинними об'єктами знаходиться за допомогою угорського алгоритму:

\begin{equation}
\hat{\sigma} = \arg\min_{\sigma \in \mathfrak{S}_N} \sum_{i=1}^N \mathcal{L}_{\text{match}}(y_i, \hat{y}_{\sigma(i)}) \tag{60}
\end{equation}

\subsubsection{Моделі для оцінки глибини: MiDaS та DPT}

Задача монокулярної оцінки глибини полягає у відновленні карти глибини за одним зображенням. Модель MiDaS \cite{ranftl2020towards} використовує багатомасштабну функцію втрат:

\begin{equation}
\mathcal{L} = \frac{1}{N} \sum_{i=1}^N \frac{1}{2\sigma_i^2} \mathcal{L}_i + \log \sigma_i \tag{34}
\end{equation}

Модель DPT (Dense Prediction Transformer) \cite{ranftl2021vision} поєднує трансформерну архітектуру з механізмом уваги:

\begin{equation}
\text{Attention}(Q, K, V) = \text{softmax}\left(\frac{QK^T}{\sqrt{d_k}}\right)V \tag{57}
\end{equation}

\subsection{Ансамблеві методи}

Ансамблеві методи дозволяють підвищити точність та надійність прогнозування шляхом комбінування кількох моделей.

\subsubsection{Метод Bagging (Bootstrap Aggregating)}

Метод Bagging передбачає усереднення результатів кількох моделей, навчених на різних підвибірках даних:

\begin{equation}
\hat{y}_{\text{bag}}(x) = \frac{1}{M} \sum_{i=1}^{M} f_i(x) \tag{53}
\end{equation}

У контексті оцінки глибини, Bagging дозволяє зменшити дисперсію прогнозу та підвищити стійкість до шумів. Експериментально показано, що використання різних аугментацій вхідних даних дає кращі результати, ніж навчання різних моделей.

\subsubsection{Метод Boosting}

На відміну від Bagging, метод Boosting будує послідовну композицію моделей, де кожна наступна модель виправляє помилки попередніх:

\begin{equation}
\hat{y}_{\text{boost}}(x) = \sum_{i=1}^{M} \alpha_i f_i(x) \tag{54}
\end{equation}

Вага кожної моделі обчислюється на основі її помилки:

\begin{equation}
\alpha_i = \frac{1}{2} \ln\left( \frac{1 - \epsilon_i}{\epsilon_i} \right) \tag{55}
\end{equation}

Для задач оцінки глибини було запропоновано адаптивне комбінування результатів різних методів:

\begin{equation}
D_{\text{final}}(x) = \begin{cases} 
D_1(x), & \text{якщо } |\nabla D_1(x)| \leq \tau \\ 
D_2(x), & \text{якщо } |\nabla D_1(x)| > \tau 
\end{cases} \tag{56}
\end{equation}

Це дозволяє використовувати переваги різних методів у різних областях зображення.

\subsection{Об'єднання методів}

Інтеграція різних методів дозволяє отримувати більш повну інформацію про сцену.

\subsubsection{Сумісна оцінка глибини та оптичного потоку (GeoNet)}

Модель GeoNet \cite{yin2019geonet} реалізує сумісну оцінку глибини та оптичного потоку на основі геометричних обмежень:

\begin{equation}
I_2(p) = I_1(p + f(p)) \tag{39}
\end{equation}

\begin{equation}
f(p) = T(p) \cdot \frac{Z(p) - Z_0}{Z(p)} \cdot \text{proj}(p) \tag{40}
\end{equation}

Такий підхід дозволяє використовувати інформацію з обох джерел для підвищення точності.

\subsubsection{Обчислення кінематичних параметрів об'єктів}

На основі отриманих полів оптичного потоку та глибини можна обчислювати кінематичні параметри об'єктів.

\textbf{Швидкість об'єкта:}

\begin{equation}
v = \frac{1}{N_\Omega} \sum_{(x,y) \in \Omega} \|F(x,y)\| \tag{41}
\end{equation}

\begin{equation}
v = \frac{1}{|\Omega|} \iint_{\Omega} \sqrt{u(x,y)^2 + v(x,y)^2} \, dx\,dy \tag{42}
\end{equation}

\textbf{Напрямок руху:}

\begin{equation}
\theta = \arctan\left( \frac{\sum_{(x,y) \in \Omega} \sin\theta(x,y)}{\sum_{(x,y) \in \Omega} \cos\theta(x,y)} \right) \tag{44}
\end{equation}

\textbf{Просторове положення:}

\begin{align}
x_c &= \frac{x_{\min} + x_{\max}}{2}, \quad y_c = \frac{y_{\min} + y_{\max}}{2} \tag{46-47} \\
z &= \text{median}\left\{ D(x,y) \mid (x,y) \in \Omega \right\} \tag{48}
\end{align}

\textbf{Траєкторія об'єкта:}

\begin{equation}
T(t) = \left\{ p(t_1), p(t_2), \ldots, p(t_n) \right\}, \quad p(t_i) = (x_i, y_i) \tag{50}
\end{equation}

\begin{equation}
\tilde{p}(t_i) = \frac{1}{2k+1} \sum_{j=-k}^{k} p(t_{i+j}) \tag{51}
\end{equation}

\subsubsection{Нормалізація та візуалізація}

Для коректного порівняння результатів різних методів використовується Min-Max нормалізація:

\begin{equation}
D_{\text{norm}}(x,y) = \frac{D(x,y) - \min(D)}{\max(D) - \min(D)} \times 255 \tag{75}
\end{equation}

Візуалізація оптичного потоку виконується шляхом перетворення у кольорове зображення:

\begin{equation}
h = \frac{\arctan(v, u)}{2\pi} \times 180^\circ \tag{79}
\end{equation}

\begin{equation}
v_{\text{bright}} = \min\left(255, \frac{\sqrt{u^2 + v^2}}{\max(\sqrt{u^2 + v^2})} \times 255\right) \tag{81}
\end{equation}

\subsection{Висновки до розділу}

У розділі 2 розглянуто основні методи дистанційної ідентифікації динамічних об'єктів на основі комп'ютерного зору та глибинного навчання.

Проведений аналіз дозволяє зробити наступні висновки:

\begin{enumerate}
    \item \textbf{Методи оптичного потоку} поділяються на локальні (Лукаса-Канаде), глобальні (Хорна-Шунка) та поліноміальні (Farneback). Кожен із методів має свої переваги: локальні — обчислювальну ефективність, глобальні — здатність відновлювати потік в однорідних областях, поліноміальні — високу точність та стійкість до шумів.
    
    \item \textbf{Нейромережеві моделі} (FlowNet, RAFT для оптичного потоку; DETR для виявлення об'єктів; MiDaS, DPT для оцінки глибини) забезпечують значно вищу точність порівняно з класичними методами, особливо в умовах складного руху та оклюзій. Використання архітектури трансформера (DETR, DPT) дозволяє моделювати довгострокові залежності в зображенні.
    
    \item \textbf{Ансамблеві методи} (Bagging, Boosting) дозволяють підвищити стійкість прогнозування шляхом комбінування кількох моделей. Особливо ефективним є адаптивне комбінування на основі локальних характеристик зображення.
    
    \item \textbf{Об'єднання методів} (GeoNet) дозволяє використовувати синергію різних джерел інформації. На основі отриманих полів оптичного потоку та глибини можна обчислювати кінематичні параметри об'єктів: швидкість (41-42), напрямок руху (44), просторове положення (46-48), траєкторію (50) та прискорення (52).
    
    \item Запропонований математичний апарат дозволяє створити інтегровану систему для дистанційної ідентифікації динамічних об'єктів, що поєднує переваги різних методів та забезпечує високу точність і надійність в реальних умовах експлуатації.
\end{enumerate}

Отримані результати є теоретичною основою для практичної реалізації системи, описаній у наступних розділах.

\begin{thebibliography}{99}

\bibitem{horn1981determining}
Horn B.K.P., Schunck B.G. (1981). Determining optical flow. \textit{Artificial Intelligence}, 17(1-3), 185-203.

\bibitem{lucas1981iterative}
Lucas B.D., Kanade T. (1981). An iterative image registration technique with an application to stereo vision. \textit{Proceedings of the 7th International Joint Conference on Artificial Intelligence}, 674-679.

\bibitem{farneback2003two}
Farnebäck G. (2003). Two-frame motion estimation based on polynomial expansion. \textit{Scandinavian Conference on Image Analysis}, 363-370.

\bibitem{dosovitskiy2015flownet}
Dosovitskiy A. et al. (2015). FlowNet: Learning optical flow with convolutional networks. \textit{Proceedings of the IEEE International Conference on Computer Vision}, 2758-2766.

\bibitem{teed2020raft}
Teed Z., Deng J. (2020). RAFT: Recurrent all-pairs field transforms for optical flow. \textit{European Conference on Computer Vision}, 402-419.

\bibitem{carion2020end}
Carion N. et al. (2020). End-to-end object detection with transformers. \textit{European Conference on Computer Vision}, 213-229.

\bibitem{ranftl2020towards}
Ranftl R. et al. (2020). Towards robust monocular depth estimation: Mixing datasets for zero-shot cross-dataset transfer. \textit{IEEE Transactions on Pattern Analysis and Machine Intelligence}, 44(3), 1623-1637.

\bibitem{ranftl2021vision}
Ranftl R. et al. (2021). Vision transformers for dense prediction. \textit{Proceedings of the IEEE/CVF International Conference on Computer Vision}, 12159-12168.

\bibitem{yin2019geonet}
Yin Z., Shi J. (2019). GeoNet: Unsupervised learning of dense depth, optical flow and camera pose. \textit{Proceedings of the IEEE/CVF Conference on Computer Vision and Pattern Recognition}, 1983-1992.

\end{thebibliography}

\end{document}